\documentclass{article}
\usepackage{amsmath,amssymb,graphicx}
\usepackage{color}
\usepackage{xcolor}
\usepackage{graphicx}
\usepackage{natbib}

\usepackage{xr-hyper}
\usepackage{hyperref}

\usepackage[margin=2.5cm]{geometry}

\setlength{\parindent}{0em}
\setlength{\parskip}{1em}

\externaldocument{../4_revised_paper/paper_sh_impacts_aerosol}

\begin{document}

%%%%%%%%%%%%%%%%%%%%%%%%%%%%%%%%%%%%%%%%%%%%%%%%%%%%%%%%%%%%%%%%%%%%%%

\section*{Response to Reviewer Report \#1}
We thank the reviewers for their careful consideration of our revised manuscript and for their insightful feedback. The page and line numbers we quote for indicating changes to the manuscript refer to the final, unmarked PDF version. 
 
%%%%%%%%%%%%%%%%%%%%%%%%%%%%%%%%%%%%%%%%%%%%%%%%%%%%%%%%%%%%%%%%%%%%%%

\textbf{(1.1)} The authors have provided highly detailed responses and conducted additional simulations addressing the points raised by the reviewers. As anticipated, simulations performed at higher relative humidity highlight the challenges in defining CCN as a function of supersaturation in cases with high concentrations of semi-volatile compounds. The partitioning of these compounds among particles of different sizes is also influenced by boundary layer dynamics. Consequently, the diagnosed CCN enhancement at different supersaturation levels is likely sensitive to the rate of relative humidity change within ascending air parcels. It is unclear whether this effect can be quantified from the simulations performed in the present study; however, the authors could expand the discussion of this aspect in the limitations section. Importantly, this is not solely a modelling issue but also a challenge for observational studies. The authors may consider including this kind of additional discussion at their discretion.

\begin{quote}
  We thank the reviewer for drawing attention to this important consideration. To address this, we have added the following statement to Section 4 (Limitations and future work) on page 18, lines 440--445:
  
  ``In the presence of substantial semivolatile material, the diagnosed CCN response may depend on how rapidly saturation is approached. Because semivolatile partitioning can evolve during humidification, CCN enhancement at a given supersaturation may be sensitive to parcel humidification history. This sensitivity is not quantified here and represents an important topic for future work. This is not solely a modeling issue; observationally-derived CCN spectra may face similar interpretational challenges when semivolatile compounds redistribute during instrument humidification or atmospheric ascent.''
\end{quote}

\end{document}
